\section{Qualitative Triplet Analysis}
\label{app:triplets}

We qualitatively analyse model--human triplet comparisons with human separation ratio $r_i(t) > 13$ for the best-performing model (\texttt{paraphrase-multilingual-mpnet}). For each comparison, we show the anchor statement and the statements judged closer to and farther from the anchor by the human rater. The separation ratio $r_i(t)$ describes how unambiguously the human rater expressed this ordering. \textbf{Correct} means that the model preserves the human ordering; \textbf{Incorrect} means that the model reverses it.

\subsection{Responsible AI --- Correct (10 of 61 shown)}
\label{app:triplets-rai-correct}

\paragraph{$r_i(t) = 15.5$}
\begin{description}[style=nextline,leftmargin=1em,font=\normalfont\bfseries]
  \item[Anchor] Kritiske kompetener vil derfor også forsvine.
  \item[Human-judged closer] Dumhed.   Vi mister personlig dømmekraft.   Faglighed er i fremtiden ikke iboende hos individer.   Vi mister kompetencer, da der ikke længere er efterspørgsel på færdigheder og kunnen.
  \item[Human-judged farther] Forskning
\end{description}
\paragraph{$r_i(t) = 15.5$}
\begin{description}[style=nextline,leftmargin=1em,font=\normalfont\bfseries]
  \item[Anchor] Eleverne og studerende lærer kun at bruge 'færdiglavede' løsninger gennem AI, i stedet for at lave sine egne.
  \item[Human-judged closer] Almen dannelse!  Manglende viden om, hvad AI kan bruges til.
  \item[Human-judged farther] Specialiseringa  Ansvarlig AI skal bidrage til at mennesker kan fokusere på at løse problemstillinger, som AI ikke kan hjælpe os med, fx menneskelig relationsarbejde.
\end{description}
\paragraph{$r_i(t) = 15.5$}
\begin{description}[style=nextline,leftmargin=1em,font=\normalfont\bfseries]
  \item[Anchor] AI tilbyder konstruerede sandheder, der kan tilpasses den enkelte behov - hvordan sikrer vi anvendelse af AI, der tilgodeser fællesskabet?
  \item[Human-judged closer] Det er - og bliver for nemt at skabe opdigtet og manipuleret indhold til de forskellige digitale platforme.
  \item[Human-judged farther] Ren hype
\end{description}
\paragraph{$r_i(t) = 15.4$}
\begin{description}[style=nextline,leftmargin=1em,font=\normalfont\bfseries]
  \item[Anchor] Bias og diskrimination  Et program er kun så godt som sin data, og det kan repræsentere et problem hvis det bliver anvendt frit ig uden regulering
  \item[Human-judged closer] AI kan karakterisere os til jobsamtaler, forsikringer og vurdere os som mennesker
  \item[Human-judged farther] Jeg går ud fra at I mener U-ansvarlig.... (slåfejl?)
\end{description}
\paragraph{$r_i(t) = 15.4$}
\begin{description}[style=nextline,leftmargin=1em,font=\normalfont\bfseries]
  \item[Anchor] Uddannelse   På uddannelser er personale ikke klædt på til at tackle AI og kan derfor ikke læse unge mennesker på til at være kritiske ifht materialer og viden skabt via AI.
  \item[Human-judged closer] Uddannelse
  \item[Human-judged farther] Ulighed    Dem som han adgang til den nyeste og kraftigste AI bliver A-holdet.
\end{description}
\paragraph{$r_i(t) = 15.4$}
\begin{description}[style=nextline,leftmargin=1em,font=\normalfont\bfseries]
  \item[Anchor] Algoritme bias  AI-systemer kan indlejre og forstærke bias
  \item[Human-judged closer] AI er ikke klogere end os i fortiden.  Jeg er bekymret for, at AI kun kommer til at forstærke stereotyper, mønstre, og uligheder i vores samfund, når det er baseret på fortiden.
  \item[Human-judged farther] Jeg går ud fra at I mener U-ansvarlig.... (slåfejl?)
\end{description}
\paragraph{$r_i(t) = 15.3$}
\begin{description}[style=nextline,leftmargin=1em,font=\normalfont\bfseries]
  \item[Anchor] AI kan karakterisere os til jobsamtaler, forsikringer og vurdere os som mennesker
  \item[Human-judged closer] Bias og diskrimination  Et program er kun så godt som sin data, og det kan repræsentere et problem hvis det bliver anvendt frit ig uden regulering
  \item[Human-judged farther] Jeg går ud fra at I mener U-ansvarlig.... (slåfejl?)
\end{description}
\paragraph{$r_i(t) = 15.2$}
\begin{description}[style=nextline,leftmargin=1em,font=\normalfont\bfseries]
  \item[Anchor] Dumhed.   Vi mister personlig dømmekraft.   Faglighed er i fremtiden ikke iboende hos individer.   Vi mister kompetencer, da der ikke længere er efterspørgsel på færdigheder og kunnen.
  \item[Human-judged closer] Kritiske kompetener vil derfor også forsvine.
  \item[Human-judged farther] Forskning
\end{description}
\paragraph{$r_i(t) = 15.0$}
\begin{description}[style=nextline,leftmargin=1em,font=\normalfont\bfseries]
  \item[Anchor] Mangel på relevant lovgivning og rel  gulering
  \item[Human-judged closer] Manglende regulering og fokus på at de skal komme alle mennesker til gode uden at ødelægge planeten. Der mangler demokratisk indflydelse fra borgernes side.
  \item[Human-judged farther] Ren hype
\end{description}
\paragraph{$r_i(t) = 15.0$}
\begin{description}[style=nextline,leftmargin=1em,font=\normalfont\bfseries]
  \item[Anchor] Dumhed.   Vi mister personlig dømmekraft.   Faglighed er i fremtiden ikke iboende hos individer.   Vi mister kompetencer, da der ikke længere er efterspørgsel på færdigheder og kunnen.
  \item[Human-judged closer] Kritiske kompetener vil derfor også forsvine.
  \item[Human-judged farther] Økonomiske interesser:  Stærke økonomiske interesser kan let stå i vejen for passende politisk regulering.
\end{description}

\subsection{Responsible AI --- Incorrect (10)}
\label{app:triplets-rai-incorrect}

\paragraph{$r_i(t) = 16.0$}
\begin{description}[style=nextline,leftmargin=1em,font=\normalfont\bfseries]
  \item[Anchor] AI tilbyder konstruerede sandheder, der kan tilpasses den enkelte behov - hvordan sikrer vi anvendelse af AI, der tilgodeser fællesskabet?
  \item[Human-judged closer] Det er - og bliver for nemt at skabe opdigtet og manipuleret indhold til de forskellige digitale platforme.
  \item[Human-judged farther] AI er ikke klogere end os i fortiden.  Jeg er bekymret for, at AI kun kommer til at forstærke stereotyper, mønstre, og uligheder i vores samfund, når det er baseret på fortiden.
\end{description}
\paragraph{$r_i(t) = 15.2$}
\begin{description}[style=nextline,leftmargin=1em,font=\normalfont\bfseries]
  \item[Anchor] Det er - og bliver for nemt at skabe opdigtet og manipuleret indhold til de forskellige digitale platforme.
  \item[Human-judged closer] AI tilbyder konstruerede sandheder, der kan tilpasses den enkelte behov - hvordan sikrer vi anvendelse af AI, der tilgodeser fællesskabet?
  \item[Human-judged farther] AI er ikke klogere end os i fortiden.  Jeg er bekymret for, at AI kun kommer til at forstærke stereotyper, mønstre, og uligheder i vores samfund, når det er baseret på fortiden.
\end{description}
\paragraph{$r_i(t) = 14.5$}
\begin{description}[style=nextline,leftmargin=1em,font=\normalfont\bfseries]
  \item[Anchor] Almen dannelse!  Manglende viden om, hvad AI kan bruges til.
  \item[Human-judged closer] Eleverne og studerende lærer kun at bruge 'færdiglavede' løsninger gennem AI, i stedet for at lave sine egne.
  \item[Human-judged farther] Specialiseringa  Ansvarlig AI skal bidrage til at mennesker kan fokusere på at løse problemstillinger, som AI ikke kan hjælpe os med, fx menneskelig relationsarbejde.
\end{description}
\paragraph{$r_i(t) = 13.8$}
\begin{description}[style=nextline,leftmargin=1em,font=\normalfont\bfseries]
  \item[Anchor] Psykiatrien
  \item[Human-judged closer] Vi har brug for et samfund med omsorg og nærvær - det kan vi ikke få digitalt
  \item[Human-judged farther] Forskning
\end{description}
\paragraph{$r_i(t) = 13.6$}
\begin{description}[style=nextline,leftmargin=1em,font=\normalfont\bfseries]
  \item[Anchor] Kritiske kompetener vil derfor også forsvine.
  \item[Human-judged closer] Vi har et skolesystem og velfærdsystem i knæ, det kan vi ikke løse digitalt
  \item[Human-judged farther] Økonomiske interesser:  Stærke økonomiske interesser kan let stå i vejen for passende politisk regulering.
\end{description}
\paragraph{$r_i(t) = 13.6$}
\begin{description}[style=nextline,leftmargin=1em,font=\normalfont\bfseries]
  \item[Anchor] AI tilbyder konstruerede sandheder, der kan tilpasses den enkelte behov - hvordan sikrer vi anvendelse af AI, der tilgodeser fællesskabet?
  \item[Human-judged closer] Det er - og bliver for nemt at skabe opdigtet og manipuleret indhold til de forskellige digitale platforme.
  \item[Human-judged farther] Hvis det bliver for restriktiv, så vil ingen eller kun få bruge det  Så mester vi en masse viden, læring og effektivisering
\end{description}
\paragraph{$r_i(t) = 13.4$}
\begin{description}[style=nextline,leftmargin=1em,font=\normalfont\bfseries]
  \item[Anchor] Dumhed.   Vi mister personlig dømmekraft.   Faglighed er i fremtiden ikke iboende hos individer.   Vi mister kompetencer, da der ikke længere er efterspørgsel på færdigheder og kunnen.
  \item[Human-judged closer] Ukontrolleret anvendelse af AI
  \item[Human-judged farther] Mener ikke det er et problem hvis der er tale om ansvarlig udvikling og anvendelse.
\end{description}
\paragraph{$r_i(t) = 13.3$}
\begin{description}[style=nextline,leftmargin=1em,font=\normalfont\bfseries]
  \item[Anchor] Facttjekke  Med al misinformationen i medierne i dag kan det blive svært at facttjekke, hvis meget indhold generet af AI og især hvis dette er bareseret på hinanden
  \item[Human-judged closer] Det er - og bliver for nemt at skabe opdigtet og manipuleret indhold til de forskellige digitale platforme.
  \item[Human-judged farther] AI er ikke klogere end os i fortiden.  Jeg er bekymret for, at AI kun kommer til at forstærke stereotyper, mønstre, og uligheder i vores samfund, når det er baseret på fortiden.
\end{description}
\paragraph{$r_i(t) = 13.3$}
\begin{description}[style=nextline,leftmargin=1em,font=\normalfont\bfseries]
  \item[Anchor] Manglende gennemsigtighed
  \item[Human-judged closer] Tillid:  Hvordan og af hvem trænes AI til at oparbejde, forstå og indeholde menneskelige egenskaber?
  \item[Human-judged farther] Manglende etiske retningslinjer og regulering  Udviklingen sker så hurtig, at lovgivning vil have svært ved at følge med
\end{description}
\paragraph{$r_i(t) = 13.1$}
\begin{description}[style=nextline,leftmargin=1em,font=\normalfont\bfseries]
  \item[Anchor] Penge  Der er rigtig mange penge involveret i at bruge data til at reklamere osv overfor danskere.
  \item[Human-judged closer] IT-giganter, vi ikke kan tøjle    Vi kan ikke styre IT-giganter som google, meta, microsoft og chatGPT. Det bliver meget svært at regulere det her område selv.
  \item[Human-judged farther] Ren hype
\end{description}

\subsection{Welfare --- Correct (3 of 39 shown)}
\label{app:triplets-welfare-correct}

\paragraph{$r_i(t) = 16.2$}
\begin{description}[style=nextline,leftmargin=1em,font=\normalfont\bfseries]
  \item[Anchor] Demografi.  Massive ændringer i befolkningssammensætningen vil betyde øgede udgifter med mange ældre med behov for pleje og sundhedsydelser.
  \item[Human-judged closer] Demografi  Meget flere ældre end unge
  \item[Human-judged farther] Øgede forsvarsudgifter
\end{description}
\paragraph{$r_i(t) = 15.6$}
\begin{description}[style=nextline,leftmargin=1em,font=\normalfont\bfseries]
  \item[Anchor] Demografi.  Massive ændringer i befolkningssammensætningen vil betyde øgede udgifter med mange ældre med behov for pleje og sundhedsydelser.
  \item[Human-judged closer] Demografi  Meget flere ældre end unge
  \item[Human-judged farther] Omkostninger til krig og oprustning
\end{description}
\paragraph{$r_i(t) = 15.5$}
\begin{description}[style=nextline,leftmargin=1em,font=\normalfont\bfseries]
  \item[Anchor] Demografi  Meget flere ældre end unge
  \item[Human-judged closer] Demografi.  Massive ændringer i befolkningssammensætningen vil betyde øgede udgifter med mange ældre med behov for pleje og sundhedsydelser.
  \item[Human-judged farther] Øgede forsvarsudgifter
\end{description}

\subsection{Welfare --- Incorrect (3)}
\label{app:triplets-welfare-incorrect}

\paragraph{$r_i(t) = 13.5$}
\begin{description}[style=nextline,leftmargin=1em,font=\normalfont\bfseries]
  \item[Anchor] Demografi.  Massive ændringer i befolkningssammensætningen vil betyde øgede udgifter med mange ældre med behov for pleje og sundhedsydelser.
  \item[Human-judged closer] Demografi  Meget flere ældre end unge
  \item[Human-judged farther] Der kommer flere og flere komplekse opgaver ud til den kommunale pleje især med ønsket om at løfte flere opgaver i det nære sundhedsområde.  Hvis ikke ressourcer er til stede, kan opgaverne blive svære at løfte
\end{description}
\paragraph{$r_i(t) = 13.4$}
\begin{description}[style=nextline,leftmargin=1em,font=\normalfont\bfseries]
  \item[Anchor] Sikring mod klimaforandringer
  \item[Human-judged closer] At grøn trepart går amok , som minkerstatningerne
  \item[Human-judged farther] Destruktive angreb udefra på vores fysiske og teknologiske infrastrukturer
\end{description}
\paragraph{$r_i(t) = 13.0$}
\begin{description}[style=nextline,leftmargin=1em,font=\normalfont\bfseries]
  \item[Anchor] For stort og tungt et administrativt system
  \item[Human-judged closer] Prioritering af andre indsatsområder (Forsvaret, skattelettelser, klima mv.) fra Christiansborg. Derfor nedprioritering af kommunernes økonomi til den borgernære velfærd.
  \item[Human-judged farther] For lille produktivitet
\end{description}
